\uuid{4rM6}
\titre{Rayon de convergence et calcul de somme }

\niveau{L2}
\module{Séries entières}
\chapitre{Séries entières}
\sousChapitre{Rayon de convergence}
\theme{Méthode du quotient}
\auteur{}
\datecreate{ 2026-04-10}
\organisation{AMSCC}
\difficulte{2}

\contenu{

\texte{ 
On considère la série entière à variable réelle
$$
T(x)=\sum_{n\ge 0}\frac{2^{n+1}}{n+1}\,x^{n+1}.
$$
}

\begin{enumerate}

\item   \question{Déterminer le rayon de convergence de la série entière $T(x)$ ainsi que son domaine ouvert de convergence.}
\indication{}
\reponse{
On pose
$$
u_n(x)=\frac{2^{n+1}}{n+1}x^{n+1}.
$$
On applique la règle de d’Alembert :
$$
\left|\frac{u_{n+1}(x)}{u_n(x)}\right|
=
\frac{2^{n+2}}{n+2}|x|^{n+2}\cdot \frac{n+1}{2^{n+1}|x|^{n+1}}
=
2|x|\frac{n+1}{n+2}.
$$
Lorsque $n\to+\infty$, on obtient
$$
\left|\frac{u_{n+1}(x)}{u_n(x)}\right|\longrightarrow 2|x|.
$$
La série converge donc si
$$
2|x|<1,
$$
c’est-à-dire si
$$
|x|<\frac12.
$$
Le rayon de convergence est donc
$$
\boxed{R=\frac12.}
$$


D’après la question précédente, la série entière converge pour
$$
|x|<\frac12.
$$
Ainsi, l’intervalle ouvert de convergence sur $\mathbb{R}$ est
$$
\boxed{\left]-\frac12,\frac12\right[.}
$$
}
\item   \question{Calculer la somme de la série sur l'intervalle ouvert de convergence.}
\indication{}
\reponse{
Pour $|x|<\frac12$, on pose
$$
y=2x.
$$
Alors $|y|<1$ et la série devient
$$
\sum_{n\ge 0}\frac{2^{n+1}}{n+1}x^{n+1}
=
\sum_{n\ge 0}\frac{(2x)^{n+1}}{n+1}
=
\sum_{n\ge 0}\frac{y^{n+1}}{n+1}.
$$
Or on sait que, pour $|y|<1$,
$$
-\ln(1+y)=\sum_{n\ge 1}\frac{(-1)^{n+1}}{n}y^n
$$
et surtout
$$
-\ln(1-y)=\sum_{n\ge 1}\frac{y^n}{n}.
$$
Donc ici,
$$
\sum_{n\ge 0}\frac{y^{n+1}}{n+1}
=
\sum_{m\ge 1}\frac{y^m}{m}
=
-\ln(1-y).
$$
En revenant à $y=2x$, on obtient
$$
\boxed{\sum_{n\ge 0}\frac{2^{n+1}}{n+1}x^{n+1}=-\ln(1-2x)}
\qquad \text{pour } |x|<\frac12.
$$
}



\end{enumerate}

}